
\documentclass[12pt]{article}

\usepackage[utf8]{inputenc}
\usepackage[T2A]{fontenc}
\usepackage[english, russian]{babel}

\usepackage[pdftex]{graphicx}
\usepackage{indentfirst}
\usepackage{float}
\usepackage{wrapfig}
\usepackage{pdflscape}
\usepackage{amssymb,amsfonts,amsmath}
\usepackage{tikz}
\usepackage{xspace}
\usepackage{changepage}
\usepackage{bytefield}
\usepackage{enumerate}

\usepackage{listingsutf8}
\lstloadlanguages{C}
\lstset
{
        language = C,
		breaklines,
		columns = fullflexible,
		flexiblecolumns,
		numbers = none,
        basicstyle = \tt\fontsize{12pt}{12pt}\selectfont,
        commentstyle = \bf,
        showtabs = false, 
        showspaces = false,
        showstringspaces = false,
        tabsize = 2,
        inputencoding = utf8/cp1251,
		frame = single,
		showlines = true,
		resetmargins = true
}

\usepackage{geometry}
\geometry{left=2cm}
\geometry{right=2cm}
\geometry{top=2cm}
\geometry{bottom=2cm}

\renewcommand{\theenumi}{\arabic{enumi}.}
\renewcommand{\theenumii}{\arabic{enumii}}
\renewcommand{\theenumiii}{.\arabic{enumiii}}
\renewcommand{\theenumiv}{.\arabic{enumiv}}

\renewcommand{\labelenumi}{\arabic{enumi}.}
\renewcommand{\labelenumii}{\arabic{enumi}.\arabic{enumii}}
%\renewcommand{\labelenumiii}{\arabic{enumi}.\arabic{enumii}.\arabic{enumiii}}
%\renewcommand{\labelenumiv}{\arabic{enumi}.\arabic{enumii}.\arabic{enumiii}.\arabic{enumiv}}

% Переопределение \caption
\makeatletter
\renewcommand{\@biblabel}[1]{#1.}
\renewcommand{\@makecaption}[2]{%
\vspace{\abovecaptionskip}%
\sbox{\@tempboxa}{#1: #2}
\ifdim \wd\@tempboxa >\hsize
   #1: #2\par
\else
   \global\@minipagefalse
   \hbox to \hsize {\hfil \small Рисунок \thefigure~<<#2>>\hfil}%
\fi
\vspace{\belowcaptionskip}}
\makeatother

\newcommand{\mysection}[1]
{
	\clearpage
	\setcounter{figure}{0}
	\renewcommand{\thefigure}{\thesection.\arabic{figure}}
	\section{#1}
}

\newcommand{\mylabsection}[1]
{
	\clearpage
	\setcounter{figure}{0}
	\renewcommand{\thefigure}{\thesection.\arabic{figure}}
	\section*{#1}
	\addcontentsline{toc}{section}{\hskip 16pt #1}
	\addtocounter{section}{1}
	\setcounter{subsection}{0}
}

\newcommand{\intro}
{
	\clearpage
	\setcounter{figure}{0}
	\renewcommand{\thefigure}{\arabic{figure}}
	\section*{Введение}
	\addcontentsline{toc}{section}{\hskip 16pt Введение}
	\input{content/intro}
}

\newcommand{\conclusion}
{
	\clearpage
	\setcounter{figure}{0}
	\renewcommand{\thefigure}{\arabic{figure}}
	\section*{Заключение}
	\addcontentsline{toc}{section}{\hskip 16pt Заключение}
	\input{content/conclusion}
}

\newcommand{\mysubsection}[1]
{
	\setcounter{figure}{0}
	\renewcommand{\thefigure}{\thesubsection.\arabic{figure}}
	\subsection{#1}
}

\newcommand{\mysubsubsection}[1]
{
	\setcounter{figure}{0}
	\renewcommand{\thefigure}{\thesubsubsection.\arabic{figure}}
	\subsubsection{#1}
}

\makeatletter
\renewcommand{\tableofcontents}
{
	\clearpage
	\par{\bf \Large \noindent \centerline{Содержание}}
	\par
	\@starttoc{toc}
}
\makeatother

\makeatletter 
\renewcommand\appendix
{
        \par
        \setcounter{section}{0}%
		\setcounter{figure}{0}%
        \gdef\thesection{\@Asbuk\c@section}
}
\makeatother

\newcommand\myappendix[1]
{
		\clearpage
        \refstepcounter{section}
		\renewcommand{\thefigure}{\thesection.\arabic{figure}}
        
        \section*{\centering{Приложение~\thesection}}
        \begin{center}
                \Large \bf #1
        \end{center}
        
        \addcontentsline{toc}{section}{Приложение~\thesection:~#1}
}

\newcommand{\mysource}[2]
{
        \refstepcounter{figure}
		\label{listing:#2}
		{
			\begin{center}

			\lstinputlisting[]{#1#2}
			\nopagebreak
			
			\medskip
			{\noindent \small Листинг~\thefigure~---~Содержимое~файла~<<#2>>}
			
			\end{center}
		}
}

% Листинги внутри локумента
%
% \mylistingbegin{Метка в listing:}{Подпись}
% \begin{lstlisting}
% \end{listing}
% \mylistingend
%

\newcommand{\mylistingcaption}{}
\newcommand{\mylistinglabel}{}

\newcommand{\mylistingbegin}[2]
{
	\refstepcounter{figure}
	\renewcommand{\mylistingcaption}{\vbox{\small \centering Листинг~\thefigure~-~#2}}
	\renewcommand{\mylistinglabel}{\label{listing:#1}}
	\begin{adjustwidth}{-\leftmargin}{\rightmargin}
	\begin{minipage}{\textwidth}
}

\newcommand{\mylistingend}
{
	\mylistingcaption
	\mylistinglabel
	\end{minipage}
	\end{adjustwidth}
}

\newcommand{\mimage}[4]
{
	\vbox
	{
		\begin{center}

			\refstepcounter{figure}
			\label{image:#1}
			
			\includegraphics[#4]{image/#2.png}

			\medskip
			{\noindent \small Рис.~\thefigure~-~#3}

		\end{center}
	}
}



\newcommand{\linux}{ОС GNU/Li\-n\-ux}
\newcommand{\virtpo}{VMWare}

\newcommand{\first}{10.0.1.128}
\newcommand{\second}{10.0.2.128}
\newcommand{\midethf}{eth0}
\newcommand{\mideths}{eth1}
\newcommand{\midf}{10.0.1.1}
\newcommand{\mids}{10.0.2.1}

\newcommand{\mytitle}{<<Курс лабораторных работ\\по дисциплине <<Вычислительные сети>>\\для ОС GNU/Linux>>}
\newcommand{\tfirst}{Настройка подключения вычислительной системы к локальной вычислительной сети}
\newcommand{\tsecond}{Протоколы транспортного уровня стека протоколов TCP/IP}
\newcommand{\tthird}{Протоколы сетевого уровня стека протоколов TCP/IP}
\newcommand{\tfourth}{Перехват сетевого трафика. Межсетевой экран netfilter}

\begin{document}

	
\begin{titlepage}

\begin{center}

Министерство образования и науки Российской Федерации \\
Рязанский Государственный Радиотехнический Университет

\bigskip

Кафедра ЭВМ

\vspace{5em}

С.И. Бабаев, М.В. Акинин

\vspace{3em}

{\Large \bf Администрирование вычислительных сетей \\ и сетевое программирование в ОС GNU/Linux}

\vspace{3em}

Методические указания\\
к лабораторным работам\\
по курсу <<Архитектура вычислительных систем и компьютерных сетей>>

\vfill

\includegraphics[width=0.5\textwidth]{image/title.png}

\vfill

Рязань \the\year

\end{center}

\thispagestyle{empty}

\end{titlepage}

\setcounter{page}{2}

\newpage

 

	\tableofcontents

	\newcommand{\labplan}[2]
	{
		\mylabsection{Лабораторная работа № #1 <<#2>>}

			\mysubsection{Цель работы}

				
Ознакомление с протоколами транспортного уровня стека протоколов TCP/IP (в частности, с протоколами TCP и UDP).

Ознакомление с принципами IP-адресации сетевых узлов в вычислительных сетях, обмен в которых осуществляется с помощью протоколов стека протоколов TCP/IPv4.

Получение навыков использования утилиты netstat для определения состояния виртуальных TCP и UDP-портов сетевых интерфейсов вычислительной системы.

Получение навыков использования утилиты ncat для организации простейшего обмена данными по протоколам TCP и UDP
между операторами удаленных друг от друга вычислительных систем.

Ознакомление с содержимых файлов /etc/services и /etc/hosts. Получение навыков использования символьных псевдонимов IP-адресов сетевых узлов и
символьных псевдонимов виртуальных портов сетевых интерфейсов.

Получение навыков сетевого программирования с использованием протоколов TCP и UDP.



			\mysubsection{Теоретическая часть}

				
\subsubsection{Общие сведения}

	\linux\ является сетевой ОС и позволяет вычислительной системе, работающей под ее управлением,
	взаимодействовать с прочими вычислительными системами в сетях самой разнообразной организации.

	Так, например, \linux\ позволяет организовать функционирование вычислительной системы в составе локальной вычислительной
	сети, построенной по одной из вариаций технологии Ethernet и использующей протоколы стека протоколов
	TCP/IP\footnote{Допускается использование как протокола IPv4, так и протокола IPv6.}
	для взаимодействия узлов сети на сетевом и транспортном уровнях сетевой модели OSI.

	В общем случае, процесс функционирования вычислительной системы, работающей под управлением \linux,
	в качестве сетевого узла локальной вычислительной сети, основанной на технологии Ethernet и стеке протоколов
	TCP/IPv4, состоит из следующих этапов:

	\begin{enumerate}

		\item Физическое подключение сетевого адаптера к материнской плате вычислительной системы;

		\item Физическое подключение сетевого адаптера к каналу связи локальной вычислительной сети;

		\item Загрузка ОС вычислительной системы;

		\item Вход в ОС оператора в качестве суперпользователя;

		\item Настройка подключения сетевого адаптера к локальной вычислительной сети - создание сетевого узла на базе вычислительной системы. 
		
		Вычислительная система обеспечивает функционирование сетевого узла, подключенного через сетевой адаптер к локальной вычислительной сети.
		Сетевой адаптер доступен процессам пользовательского уровня ОС через соответствующий сетевой интерфейс.

		Данный этап предполагает:

			\begin{itemize}

				\item Задание IPv4-адреса настраиваемого сетевого узла;
				\item Задание маски сети\footnote{Маска сети суть есть целое беззнаковое число, позволяющее определить адрес сети и широковещательный адрес
				сети по адресу любого сетевого узла, состоящего в сети.} и широковещательного адреса сети\footnote{Широковещательный адрес сети используется
				при сетевом обмене для массовой рассылки пакетов всем сетевым узлам, состоящим в сети.}, в которую включен настраиваемый сетевой узел;
				\item Настройка дополнительных параметров сетевого узла.
				
				К числу дополнительных параметров сетевого узла относится, например, Maximum Transmission Unit
				(MTU) - максимальный размер пакета, могущего быть переданным через данный сетевым узлом или данному сетевому узлу
				по протоколу канального уровня сетевой модели OSI;

			\end{itemize}

		\item Настройка таблицы маршрутизации.

		С помощью таблицы маршрутизации вычислительная система принимает решение, по какому маршруту
		(какому сетевому узлу через какой сетевой интерфейс) необходимо передать очередной пакет данных, чтобы
		означенный пакет достиг сетевого узла - приемника;

		\item Запуск и настройка дополнительных сервисов.

		В число дополнительных сервисов входят, например, DNS-сервер, DHCP-сервер, WEB-сервер, FTP-сервер и прочие сервисы;

		\item Обмен данными между процессами вычислительной системы и процессами удаленных вычислительных систем;

		\item Останов дополнительных сервисов;

		\item Очистка таблицы маршрутизации;

		\item Сброс настроек сетевого интерфейса;

		\item Останов ОС.

	\end{enumerate}

	В данной лабораторной работе будет рассмотрен способ подключения вычислительной системы, работающей
	под управлением \linux, к локальной вычислительной сети, основанной на технологии Ethernet и
	стеке протоколов TCP/IPv4.

\subsubsection{Настройка подключения вычислительной системы к локальной вычислительной сети}

	Настройка подключения вычислительной системы к локальной вычислительной сети может быть произведена
	различными способами. Каждый из способов предполагает определенную степень автоматизации процесса подключения -
	от полностью ручного способа, рассматриваемого в данной лабораторной работе, до полностью автоматического,
	заключающегося в запуске операционной системой определенного набора программ с определенными ключами по
	составленному администратором ОС скрипту на некотором языке программирования\footnote{Чаще всего в качестве языка автоматизации
	выбирается язык командного интерпретатора, используемого администратором.}.

	Ручной способ настройки подключения вычислительной системы к локальной вычислительной сети предполагает следующую последовательность
	действий:

	\begin{enumerate}

		\item Вход оператора ОС в ОС в качестве суперпользователя;

		\item Настройка ядра ОС.

		Настройка ядра ОС предполагает загрузку в пространство ядра кода сетевой подсистемы ядра
		и драйвера используемого сетевого адаптера.
		
		Как правило, хотя и необязательно, код сетевой подсистемы ядра находится в файле ядра и,
		соответственно, загружается в ядерное пространство непосредственно в ходе загрузки ядра ОС на
		ранней стадии загрузки ОС. В том случае, если код сетевой подсистемы ядра скомпилирован в виде отдельных
		модулей ядра, подгружаемых в ходе работы ОС, данные модули загружаются в ходе запуска графической подсистемы
		ОС, поскольку она не может функционировать без сетевой подсистемы, или загружаются ОС автоматически,
		при выполнении попыток настройки сетевого адаптера. В крайне редких случаях, когда код сетевой подсистемы
		отсутствует в ядре, требуется перекомпилировать ядро ОС, что выходит за рамки данной лабораторной работы.

		Драйвер используемого сетевого адаптера может находится в файле ядра ОС, в отдельном модуле ядра,
		загружаемым ядром при выполнении попытки настройки сетевого адаптера, или отсутствовать в ядре вовсе.

		Определить, присутствует ли драйвер сетевого адаптера в пространстве ядра ОС, можно с помощью утилиты
		ifconfig - для этого необходимо запустить данную утилиту с ключом <<-a>>, который предписывает утилите
		ifconfig вывести в стандартный поток вывода информацию о всех присутствующих в вычислительной системе
		(но, возможно, ненастроенных) сетевых адаптерах.
		На рисунке \ref{image:lab1-5} приведен возможный вывод утилиты ifconfig, запущенной с ключом <<-a>>.

		\mimage{lab1-5}{1/5}{Список сетевых адаптеров, присутствующих в вычислительной системе}{}

		Как видно из рисунка \ref{image:lab1-5} к вычислительной системе подключены следующие сетевые адаптеры:

		\begin{itemize}

			\item eth0 - Ethernet-адаптер. IP-адрес отсутствует => адаптер не настроен;
			\item lo - виртуальный интерфейс обратной петли;
			\item vboxnet0 - виртуальный Ethernet-адаптер, используемый виртуальными машинами VirtualBox
				  (ни одна из таковых не запущена => vboxnet0 не настроен (IP-адрес отсутствует));
			\item vmnet1, vmnet8 - виртуальные Ethernet-адаптеры, используемые виртуальными машинами VMWare.
			Оба адаптера настроены и имеют следующие IP-адреса:

				\begin{itemize}

					\item vmnet1 - 192.168.126.1;
					\item vmnet8 - 192.168.91.1;

				\end{itemize}

			\item wlan0 - Wi-Fi карта (не настроена);

		\end{itemize}

		\item Настройка подключения сетевого интерфейса сетевого адаптера, с помощью которого выполняется подключение к локальной
		вычислительной сети, к локальной вычислительной сети.

		Означенная настройка предполагает:

		\begin{itemize}

			\item Задание IP-адреса настраиваемого сетевого узла;
			\item Задание маски сети и широковещательного адреса сети, в которую включен сетевой узел;
			\item Настройка дополнительных параметров сетевого узла.

		\end{itemize}

		Настройка сетевого интерфейса выполняется с помощью утилиты ifconfig и предполагает следующую последовательность действий.

		\begin{enumerate}

			\item Определение имени сетевого интерфейса.

			Интерфейсы Ethernet-адаптеров именуются в \linux\ в виде последовательности символов <<ethN>>, где N -
			порядковый номер сетевого адаптера (считая с 0), зависящий от его физического подключения к
			материнской плате вычислительной системы.

			Узнать имя сетевого интерфейса можно с помощью утилиты ifconfig, запущенной с ключом <<-a>>;

			\item Сброс настроек сетевого интерфейса.

			Для сброса настроек сетевого интерфейса необходимо выполнить команду <<ifconfig ethN down>>;

			\item Настройка сетевого интерфейса.

			Для настройки сетевого интерфейса необходимо выполнить команду <<ifconfig ethN IP broadcast BROADCAST netmask NETMASK>>, где:
			
			\begin{itemize}
			
				\item IP - IP-адрес сетевого узла;
				\item BROADCAST - широковещательный адрес сети;
				\item NETMASK - маска сети.

			\end{itemize}

			Передав утилите ifconfig дополнительные ключи, можно произвести настройку прочих параметров
			сетевого адаптера. Так, например, с помощью ключа <<mtu MTU\_SIZE>> оператор может настроить
			максимальный размер пакета (MTU\_SIZE; указывается в байтах),
			могущего быть переданным через сетевой адаптер по протоколу канального уровня сетевой модели OSI,
			что полезно для некоторых сетевых адаптеров отдельных производителей, требующих для своего корректного функционирования
			значение MTU меньшее, чем традиционное значение MTU, равное 1500 байтам.

		\end{enumerate}

		В случае успешного завершения процесса настройки сетевого интерфейса, запись о данном сетевом интерфейсе появится
		в списке корректно настроенных сетевых интерфейсов, могущего быть полученным с помощью утилиты ifconfig, запущенной без ключей,
		что проиллюстрировано рисунком \ref{image:lab1-4};

		\mimage{lab1-4}{1/4}{Список корректно настроенных сетевых интерфейсов вычислительной системы}{width=\textwidth}

		\item Настройка таблицы маршрутизации.

		После осуществления настройки сетевого интерфейса необходимо выполнить настройку таблицы маршрутизации, для чего используется утилита route.

		Одна запись, связанная с настраиваемым сетевым интерфейсом, в таблице маршрутизации уже есть -
		это указание направлять все пакеты, следующие в соответствующую сеть, через настроенный сетевой интерфейс.

		Необходимо добавить в таблицу маршрутизации запись о маршруте по умолчанию -
		то есть запись о маршруте, по которому будут отправляться пакеты в случае, если других маршрутов,
		по которым можно было бы отправить пакет, нет.
		
		Для добавления записи в таблицу маршрутизации о маршруте по умолчанию необходимо выполнить следующие действия:

		\begin{enumerate}

			\item Удалить запись о маршруте по умолчанию из таблицы маршрутизации.

			Для удаления записи о маршруте по умолчанию из таблицы маршрутизации необходимо выполнить команду
			<<route del default>>;
			
			\item Добавить запись о маршруте по умолчанию в таблицу маршрутизации
			с предписанием использовать целевой сетевой узел в качестве шлюза между сетями.

			Для добавления записи о маршруте по умолчанию в таблицу маршрутизации необходимо выполнить команду
			<<route add default gw IP ethN>>, где IP - IP-адрес сетевого узла - шлюза.

		\end{enumerate}

		Содержимое таблицы маршрутизации можно получить с помощью утилиты route, запустив данную утилиту без каких бы то ни было ключей.
		Пример получения содержимого таблицы маршрутизации с помощью утилиты route приведен на рисунке \ref{image:lab1-route}.

		\mimage{lab1-route}{1/route}{Состояние таблицы маршрутизации}{width=\textwidth}

	\end{enumerate}

\subsubsection{Проверка корректности функционирования локальной вычислительной сети}

	По завершению настройки вычислительной сети разумно проверить корректность ее функционирования,
	для чего можно использовать утилиту ping.
		
	Утилита ping предназначена для проверки доступности целевого сетевого узла и позволяет оценить
	некоторые параметры процесса обмена данными с целевым удаленным сетевым узлом.

	Для проверки доступности сетевого узла, объединенного с сетевым узлом, на котором производится
	запуск утилиты ping, в вычислительную сеть, утилите ping необходимо передать IP-адрес сетевого узла,
	доступность которого проверяется.

	На рисунке \ref{image:lab1-ping} приведен результат выполнения утилиты ping,
	запущенной на некотором сетевом узле последовательно с ключами <<-c 5 10.0.1.128>> и
	<<-c 5 10.0.2.128>>\footnote{Ключ <<-c 5>> суть есть указание использовать при проверки доступности
	целевого сетевого узла только 5 циклов обмена пакетами типов ECHO и ECHO-REPLY протокола ICMP.}.
	Как видно из рисунка \ref{image:lab1-ping}, сетевые узлы 10.0.1.128 и 10.0.2.128 доступны сетевому узлу, на котором запущена утилита ping,
	что позволяет говорить о корректности настроек локальной вычислительной сети.

	\mimage{lab1-ping}{1/ping}{Проверка корректности функционирования вычислительной сети}{width=\textwidth}



			\mysubsection{Практическая часть}

				
\subsubsection{Предварительная подготовка}

Для выполнения данной лабораторной работы должна быть использована вычислительная сеть, созданная и настроенная в ходе выполнения лабораторной работы № 1.
Структурная схема используемой вычислительной сети приведена на рисунке \ref{image:struct3}.

\mimage{struct3}{1/struct}{Структурная схема вычислительной сети}{width=\textwidth}

\subsubsection{Сетевое программирование с использованием протоколов ICMP и ARP. Утилиты ping и arp. Файл /proc/net/arp}
\label{task:l3t1}

	В рамках данной лабораторной работы каждому студенту группы необходимо выполнить на центральном сетевом узле
	следующие задания:

	\begin{enumerate}

		\item Проверить доступность каждого из периферийных сетевых узлов с помощью N циклов обмена с периферийными
		сетевыми узлами пакетами типа ECHO и ECHO-REPLY протокола ICMP;

		\item Получить содержимое ARP-кэша с помощью файла /proc/net/arp;

		\item Удалить из ARP-кэша записи об обоих периферийных сетевых узлах;

		\item Проверить корректность удаления записей из ARP-кэша с помощью утилиты arp;

		\item Разработать программу, выполняющую ARP-запрос о периферийном сетевом узле \linebreak \first;\label{lab3:pp-arp}

		\item Разработать программу, выполняющую обмен пакетами типов ECHO и ECHO-REPLY протокола ICMP с периферийным
		узлом \second.\label{lab3:pp-icmp}
		
		Программа должна выполнять последовательно N циклов обмена пакетами означенного типа с
		периферийным узлом \second\ и для каждого принятого или отправленного пакета программа должна выводить
		содержимое заголовка пакета;

		\item Запустить на выполнение программы, разработанные в пунктах \ref{lab3:pp-arp} и \ref{lab3:pp-icmp};

		\item По завершению выполнения программ, разработанных в пунктах \ref{lab3:pp-arp} и \ref{lab3:pp-icmp},
		получить содержимое ARP-кэша с помощью файла /proc/net/arp;

		\item Получить содержимое ARP-кэша с помощью утилиты arp;

		\item Сделать выводы о принципах функционирования ARP-кэша.

	\end{enumerate}

	Здесь N суть есть: $N = (M~mod~5 + 1) * 4$, где M - порядковый номер студента в журнале группы.

\subsubsection{Подготовка отчета по лабораторной работе}

	Подготовить отчет, в котором необходимо привести подробное описание процесса выполнения заданий пункта
	\ref{task:l3t1}, проиллюстрированное достаточным количеством снимков экрана.

	В отчете по лабораторной работе должен быть приведен исходный код разработанного программного обеспечения с
	соответствующими комментариями.



			\mysubsection{Защита отчета по лабораторной работе}

				
Одно из следующих заданий может быть предложено в качестве задания для защиты отчета по лабораторной работе:

\begin{enumerate}

	\item Переработка серверной компоненты, осуществлявшей обмен информацией с управляющей компонентой по протоколу
	UDP, для использования при означенном обмене протокола TCP;

	\item Переработка серверной компоненты, осуществлявшей обмен информацией с управляющей компонентой по протоколу
	TCP, для использования при означенном обмене протокола UDP;

	\item Добавление функционала обработки управляющей компонентой одной из следующих команд:

		\begin{itemize}

			\item <<get NODE FILE\_IN FILE\_OUT>> - команда запроса у серверной компоненты, запущенной на периферийном
			сетевом узле с номером NODE, содержимого файла FILE\_IN и сохранение полученной информации на центральном
			сетевом узле в файле FILE\_OUT;

			\item <<put NODE FILE\_IN FILE\_OUT>> - копирование файла FILE\_IN центрального сетевого узла в файл FILE\_OUT
			периферийного сетевого узла с номером NODE;

			\item <<route NODE>> - запрос у серверной компоненты, функционирующей на периферийном сетевом узле с номером
			NODE, таблицы маршрутизации данного периферийного сетевого узла;

		\end{itemize}

	\item Добавление поддержки обмена сообщениями между оператором управляющей компоненты и операторами сетевых узлов,
	на которых запущены серверные компоненты;

	\item Добавление поддержки симметричного шифрования данных, пересылаемых между серверными и управляющей компонентами.

		В качестве простейшего алгоритма симметричного шифрования можно использовать <<затенение>> байт с помощью
		операции побитового исключающего или.

		Ключи шифрования должны быть сгенерированы управляющей компонентой и разосланы серверным компонентам
		непосредственно после установа соединения между серверными и управляющей компонентами.

\end{enumerate}


	}

	\labplan{1}{\tfirst}
	\labplan{2}{\tsecond}
	\labplan{3}{\tthird}
	\labplan{4}{\tfourth}

	\appendix

	\myappendix{Некоторые приемы работы с ОС GNU/Linux с использованием интерфейса командной строки}

		
Пользователи, существующие в установках ОС GNU/Linux, используемых в настоящем курсе лабораторных работ:

\begin{itemize}

	\item Суперпользователь root (пароль - toor);
	\item Пользователь user (пароль - user).

\end{itemize}

ОС GNU/Linux позволяет пользователю работать с несколькими независимыми друг от друга виртуальными консольными и графическими терминалами. Переключение между консольными виртуальными терминалами осуществляется с помощью сочетаний клавиш <<Alt + FНОМЕР>> и <<Ctrl + Alt + FНОМЕР>>, где <<НОМЕР>> - номер виртуального терминала, на который оператор ОС желает переключится. Так, например, нажатие клавиш <<Alt + F1>> позволяет переключится на первый консольный виртуальный терминал, <<Alt + F2>> - на второй и <<Alt + F3>> - на третий.

Следующие команды могут оказаться полезными при работе с ОС GNU/Linux:

\begin{itemize}

	\item <<telinit 0>> - останов ОС, выключение компьютера (необходимы права суперпользователя);
	\item <<telinit 6>> - перезагрузка ОС (необходимы права суперпользователя);
	\item <<gcc ФАЙЛ.c -o ФАЙЛ\_2>> - компиляция программы, написанной на языке программирования C, исходный код которой сохранен в файле ФАЙЛ.c. На выходе генерируется исполняемый файл ФАЙЛ\_2.
	
		Для запуска на выполнение исполняемого файла ФАЙЛ\_2 оператор должен выполнить команду <<./ФАЙЛ\_2>>;

	\item <<mc>> - запуск файлового менеджера <<Midnight Commander>>, обладающего, кроме всего прочего, удобным текстовым редактором с подсветкой синтаксиса различных языков программирования;
	\item <<nano>> - запуск текстового редактора Nano;
	\item <<su ->> - вход в ОС в качестве суперпользователя;
	\item <<man СТРАНИЦА>> - получение справочной информации (содержимое man-страницы <<СТРАНИЦА>>; например, команда <<man printf>> позволяет получить справочную информацию о функции printf() стандартной библиотеки языка программирования C).
	
		Для завершения выполнения утилиты man необходимо нажать клавишу <<q>>;

	\item <<ps -A>> - получить список всех процессов системы, существующих в ОС на момент вызова утилиты <<ps>>;

	\item <<kill PID>> или <<killall ИМЯ\_ИСПОЛНЯЕМОГО\_ФАЙЛА>> - останов процесса, обладающего указанным идентификатором (PID) или соответствующего указанному исполняемому файлу.

		Другой способ останова некоторого процесса ОС - нажатие сочетания клавиш <<Ctrl + C>> в том виртуальном терминале, в котором процесс запущен. Нажатие указанного сочетания клавиш позволяет останавливать только те процессы ОС, которые были запущены на <<переднем плане>> (то есть не в фоновом режиме); 

	\item <<net\_help>> - получение дополнительной справочной информации.

		Команда <<net\_help>> уникальна для установок ОС GNU/Linux, используемых в данном курсе лабораторных работ.

\end{itemize}

При запуске исполняемых файлов на выполнение в конце командной строки не стоит указывать символ <<\&>> - если означенный символ будет указан, то процесс, соответствующий исполняемому файлу, будет запущен в фоновом режиме и, таким образом, будет, в рамках данного курса лабораторных работ, бесполезен.



\end{document}

